\documentclass[10pt,a4paper]{amsart}
\usepackage[margin=19mm]{geometry}
\usepackage{amsmath,amssymb,mathtools}
\usepackage[scr=rsfs,cal=euler]{mathalfa}
\usepackage{xfrac}
\usepackage[unicode,hidelinks,bookmarks=false]{hyperref}
\newcommand{\ie}{i.e.\ }
\newcommand{\cf}{cf.\ }
\newcommand{\resp}{respectively\ }
\newcommand{\Subsection}{\subsection}
\providecommand{\nobreakdash}{\nobreak\hbox{-}}
\setlength{\emergencystretch}{2em}
\multlinegap0pt
\theoremstyle{plain}
\newtheorem{theorem}{Theorem}
\newtheorem{lemma}{Lemma}
\newtheorem{corollary}{Corollary}
\newtheorem{definition}{Definition}
\theoremstyle{remark}
\newtheorem{remark}{Remark}
\renewcommand{\qedsymbol}{$\blacksquare$}
\title[The Kalman--Bucy filter as a linear function of the signature]{Nowcasting using regression on signatures: the continuous time Kalman--Bucy filter as a linear function of the signature of the augmented observation process}
\author{Samuel N. Cohen, Giulia Mantoan, Lars Nesheim, \'{A}ureo de Paula, Arthur Turrell, Lingyi Yang}
\begin{document}
\maketitle
\noindent\emph{Source and licence.} Nowcasting using regression on signatures. Version arXiv:2305.10256v3 (CC BY 4.0), distributed under the Creative Commons Attribution 4.0 International licence, \url{http://creativecommons.org/licenses/by/4.0/}, as recorded in the arXiv licence metadata for this exact version and shown on the abstract page \url{https://arxiv.org/abs/2305.10256v3}. Full licence notice: licenses/arxiv-2305.10256-CC-BY-4.0.txt.\par\smallskip

\noindent\textit{Editorial scope.} Counted unit: the article's Theorem with source label \texttt{thm:KalmanRepresentation} (source0.tex lines 371--373), delivered together with the article's complete original proof of it (source0.tex lines 376--405, one \texttt{proof} environment of 30 lines, adjacent to the statement). No mathematics is written, repaired or re-derived: statement and proof are the article's own text line for line, in the article's own notation, and no retained line is substituted, re-flowed or omitted. Required context retained uncounted, in source order: the article's section heading \texttt{Theory} (line 199) and the labelled subsection heading \texttt{Continuous time state-space models} (line 238), which the counted proof cites, so that the reference resolves inside the delivered document; the article's Kalman--Bucy filter equations together with the definitions of the filter estimate and of the conditional error covariance that immediately follow them (lines 249--257), which the counted statement and proof cite; and the definition of the finite one-variation norm together with the factorial bound on the magnitude of signature terms (lines 339--343), which the counted proof cites. Nothing else of the article is retained: its other sections, its simulations, its empirical applications, its appendices, its figures and tables are omitted. No citation command occurs in any retained line, so the article's bibliography is not delivered and the wrapper carries no bibliography entry. Every cross-reference inside the retained spans resolves inside the delivered document, and the labels defined inside the counted proof itself are available there. Editorial formatting only: an AMS \texttt{amsart} preamble that declares the article's theorem-like environments and restates the end-of-proof symbol of the article; the wrapper adds no mathematical macro, because every control sequence used by the retained lines is standard. No class or style file is shipped with this package, and the article's own \texttt{article}-class format with \texttt{authblk} affiliations is not redistributed; the wrapper is an amsart document. Numbering differs from the article, because the wrapper retains only the headings listed above and no section reset: the retained section prints as Section 1, the retained labelled subsection prints as Section 1.1, and the article's equations print with the wrapper's continuous numbering, so the references in the retained text point at the retained displays rather than at the article's printed numbers. No figure environment and no graphics-inclusion command occurs in any retained line, and the package delivers no image asset. One bundled source file, \texttt{sources/141650/source0.tex}, accompanies the delivered item as the exact source of every retained span. Every retained span is recorded in \texttt{delivery-evidence.json} with method \texttt{exact\_tex}.
\section{Theory}\label{sec:background}

\subsection{Continuous time state-space models}\label{sec:cts_kf}

\begin{align}
    d\hat Y_t &= (F \hat Y_t +f)dt + R_tH^\top (dX_t -(H\hat Y_t+h)dt),\label{eqn:opt_filter}\\
    \frac{dR_t}{dt} &= \sigma \sigma^\top+F R_t + R_t F^\top - R_t H^\top H R_t. \label{eqn:filter_var}
\end{align}
Here, \(\hat{Y}_t = \mathbb{E}[Y_t \mid \mathcal{F}^X_t]\) denotes the optimal estimate (in the mean square sense) of the unobserved state \(Y_t\), conditioned on the information available from the observed process \(X\) up to time \(t\), that is, \(\mathcal{F}^X_t = \sigma(X_s : s \leq t)\).
The matrix \(R_t \in \mathbb{R}^{d \times d}\) is the conditional covariance of the estimation error
\[
R_t = \mathbb{E}[(Y_t - \hat{Y}_t)(Y_t - \hat{Y}_t)^\top \mid \mathcal{F}^X_t].
\]

\[||X||_{1,[0,T]} \equiv \sup_{P \in \mathcal{P}} \sum_{i=0}^{n_P-1} |X_{i+1} - X_i| <\infty,\]
where \(\mathcal{P}\) is the set of all partitions of \([0,T]\) and \(n_P\) is the number of points in partition \(P\), then, for each level \(k \geq 1\), 
\begin{equation}\label{eqn:sig_magnitude}
    S^k(X\dots X)_{[0,T]} = \int_0^{T}\cdots\int_0^{t_{1}} dX_{t_0}  \dots dX_{t_{k-1}} \leq \frac{||X||^k_{1,[0,T]}}{k!}.
\end{equation}

\begin{theorem}\label{thm:KalmanRepresentation}
The continuous time Kalman--Bucy filter $\hat Y$ (including the case with deterministic time-varying parameters), as described by Equation \eqref{eqn:opt_filter}, can be written as a linear function of the initial estimate $\hat Y_0$ and the signature of the augmented observation process $(t,X_t)$. In particular, this representation uses only the signature terms where the observation $X$ appears once in the iterated integral.
\end{theorem}

\begin{proof} With the natural time-varying change to the notation set out in Section~\ref{sec:cts_kf}, write
\begin{align}
A_t &\equiv F_t -R_t H_t^\top H_t,\nonumber\\
\xi_t &\equiv  \int_0^t (f_s - R_s H_s^\top h_s) ds + \int_0^t R_sH_s^\top dX_s, \label{eqn:xi}
 \end{align}
which gives us the simplified expression for the filter
\begin{equation}\label{eq:simplifiedXhat} d\hat Y_t = A_t \hat Y_tdt + d\xi_t.\end{equation}

Denote the iterated integrals of $A$ and $\xi$ by
\begin{align*}
    \mathbb{A}^n_t &= \int_0^t\int_0^{t_1}\cdots\int_0^{t_{n-1}}\Big(A_{t_1}A_{t_2}\cdots A_{t_n}\Big)dt_{n}\cdots dt_1,\\
    \Xi^n_t &= \int_0^t\int_0^{t_1}\cdots\int_0^{t_{n-1}}\Big(A_{t_1}A_{t_2}\cdots A_{t_n}\xi_{t_n}\Big) dt_{n}\cdots dt_1,
\end{align*}
with the conventions $\mathbb{A}^0 =\mathrm{I}_d$ (the identity matrix) and $\Xi^0_t = \xi_t$. Observe that these are both formally solutions of the recurrence relation $Q^n_t = \int_0^t A_s Q^{n-1}_s ds$ with different values for $Q^0$ (and in different dimensions). We note that the $n$th iterated integral is super-polynomially small (that is, it satisfies Equation (\ref{eqn:sig_magnitude})), and in particular the infinite sums $\sum_{n\ge 0} \mathbb{A}^n_t$ and $\sum_{n\ge 0} \Xi^n_t$ are both well defined.


Writing $\hat Y$ in integral form, we have
\[\hat Y_t = \hat Y_0 +\int_0^t A_s \hat Y_s ds + \xi_t.\]
This integral equation has a unique solution, which can be represented as the series
\[\hat Y_t = \sum_{n\ge 0} \Big(\mathbb{A}^n_t \hat Y_0 + \Xi^n_t\Big).\]
Given the recurrence relation satisfied by $\mathbb{A}^n_t$ and $\Xi^n_t$, we can see that this is a solution to the integral equation.
This shows that $\hat Y_t$ is a (linear) function of its initial value $\hat Y_0$ and the iterated integral processes $\mathbb{A}^n$ and $\Xi^n$. This expansion is very closely related to the Picard series approximation of the stochastic differential equation defining $\hat Y$ (Equation \eqref{eqn:opt_filter}). 

It remains to show that $\mathbb{A}^n$ and $\Xi^n$ can be expressed in terms of the (joint) signature of the time-augmented path $(t, X_t)$. As $A$ is a continuous deterministic function of time, we know that (over any finite time horizon) it can be approximated arbitrarily well by a polynomial (by the Stone--Weierstrass theorem). As the signature of time is simply the sequence $\mathbb{S}(t)= (1, t, t^2/2, \dots, t^n/n!,\dots)$, we see that $A$ can be written as a (matrix-valued) linear function of the signature of $t$.

The process $\xi$ in Equation~(\ref{eqn:xi}) is slightly more delicate, as it depends on both time and the observations $X$. The first term $\int_0^t (f_s - R_s H_s^\top h_s) ds$ is deterministic, so again has a polynomial expansion in terms of the signature of $t$. Considering the second term $\int_0^t R_sH_s^\top dX_s$, we see that if $H$ is continuous\footnote{If $H$ is not continuous, then we need to use the fact that polynomials in time are dense in the $L^2([0,t])$ space, which is the appropriate space to consider given we have an outer integral with respect to the process $X$. In this case we will still have a polynomial approximation $p_n(s) \approx R_sH_s^\top$, such that the integral $\int_0^t p_n(s) dX_s \to \int_0^t R_s H_s^\top dX_s$ converges in a mean-square sense as $n\to \infty$.}, $R_sH_s^\top$ has an expansion in terms of the signature of time, so $\int_0^t R_sH_s^\top dX_s$ has an expansion in terms of the signature of $(t,X_t)$, with the special form where the only integral with respect to $X$ is the outermost one.

As both $A$ and $\xi$ have expansions in terms of the signatures of $t$ and $(t,X)$ respectively, it follows that their iterated integrals $\mathbb{A}$ and $\Xi$ also have expansions of this type. For our purposes the explicit values of this expansion are not of particular interest (and do not have a simple algebraic form), but the definition of $\mathbb{A}^n$ immediately shows that if $A$ can be written as a polynomial in time, then so can $\mathbb{A}^n$ for each $n$. Similarly for $\Xi^n$, but now this will involve iterated integrals with a single $X$ integral included.

\end{proof}

\end{document}
